\documentclass[11pt,a4paper]{article}
\usepackage[utf8]{inputenc}
\usepackage[T1]{fontenc}
\usepackage[french]{babel}
\usepackage[margin=2.2cm]{geometry}
\usepackage{amsmath,amssymb}
\usepackage{booktabs}
\usepackage{hyperref}
\usepackage{enumitem}

\title{Experiments v3 : combinaison de contraintes structurelles et appariement pent--hept pour la planarite}
\author{Stage M2 IMD --- Generateur non-benzenoides}
\date{Mai 2026}

\begin{document}
\maketitle

\section{Motivation et continuite avec v2}

\texttt{experiments\_v2} a identifie la contrainte \texttt{pb1}
($K_{\text{pb}} = 1$, au plus un pentagone sur le bord du squelette)
comme le levier dominant pour ameliorer le taux de planarite des
molecules generees, avec un gain de $+9$ a $+19$ points selon
$h \in \{6, 7, 8, 9\}$. La symetrie globale 5/7 (\texttt{sym1}) ne
produit aucun gain mesurable sur les configurations testees, et le
durcissement \texttt{pb0} rend le probleme insatisfiable.

L'objectif d'\texttt{experiments\_v3} est de pousser le levier
combinatoire plus loin, en exploitant une propriete \emph{geometrique
locale} qui n'etait pas exprimee dans v2 : l'\emph{appariement pent--hept
de proximite}. L'intuition vient de la theorie du \emph{defaut
topologique} en physique des materiaux 2D (defaut Stone--Wales) : un
pentagone isole introduit une courbure positive locale ($+\pi/3$ par
Gauss--Bonnet discret) qui ne peut etre compensee que par un heptagone
spatialement proche (courbure $-\pi/3$). Si chaque pentagone est apparie
a un heptagone dans son voisinage immediat, la courbure cumulee
localement est nulle, condition necessaire pour que la structure puisse
se plonger dans un plan.

\section{Nouvelle contrainte : combinaison pb1 + adj-57}

La contrainte \texttt{adj-57} existait deja dans le solveur (drapeau
optionnel) mais n'avait pas ete combinee a \texttt{pb1}. Sa
formalisation est :
\begin{align*}
  \text{(C5)} \quad & \forall v,\ x_v = 5 \;\Rightarrow\; \sum_{u \in N(v)} \mathbf{1}[x_u = 7] \ge 1 \\
                    & \forall v,\ x_v = 7 \;\Rightarrow\; \sum_{u \in N(v)} \mathbf{1}[x_u = 5] \ge 1
\end{align*}
ou $N(v)$ designe les voisins de $v$ dans le graphe dual. Elle impose
que chaque pentagone ait au moins un heptagone parmi ses voisins
immediats, et reciproquement. Combinee a \texttt{pb1}, qui contraint la
localisation globale des pentagones, on obtient une configuration ou la
courbure positive (au plus un pent en bord) est non seulement
geographiquement limitee, mais doit aussi etre compensee \emph{a une
distance d'un seul cycle} par une courbure negative. C'est l'analogue
discret du \emph{quadripole Stone--Wales} $5|7|7|5$ qui caracterise
les defauts plans du graphene.

Cette combinaison est implementee comme un nouveau preset
\texttt{pb1\_adj57} dans \texttt{experiments\_v2/configs.py} (memes
valeurs $K_{\text{pb}} = 1$ que \texttt{pb1}, le drapeau \texttt{adj-57}
etant passe via le mecanisme \texttt{extra-flag} pour distinguer les
deux experiences dans la base de donnees finale \texttt{db\_v3.db}).

\section{Pipeline et une note sur les filtres post-CSP}

Le pipeline d'execution reprend strictement l'architecture v2 : un
manifest JSONL liste les couples $(\text{squelette}, \text{config})$ a
traiter, chaque job lance la generation CSP, reconstruit les molecules
en 3D et valide la planarite par xTB MD + optimisation finale. La base
\texttt{db\_v3.db} agrege l'integralite des resultats par couple
$(h, \text{config})$, permettant la comparaison directe avec
\texttt{pb1}, \texttt{pb2} et les configurations de v2.

Une digression a ete tentee en parallele : ajouter un \emph{pre-filtre
geometrique rapide post-CSP} (MMFF94 par RDKit) afin de ne lancer xTB
que sur les candidats que MMFF ne sait pas immediatement aplatir. Sur
une validation initiale en aveugle sur des geometries deja optimisees
par xTB (extraites de \texttt{db\_v2}), MMFF94 atteignait une accuracy
de 80--84~\%, ce qui semblait acceptable. Mais une validation par
echantillonnage independante (5\,996 molecules tirees au sort dans le
sous-ensemble \emph{mmff sure plan} d'un run reel, relancees par xTB
complet) a revele une precision reelle de seulement 58~\% : MMFF94
fonctionne sur des geometries deja a l'equilibre, mais echoue
massivement sur les geometries d'entree brutes generees par la
reconstruction (effets electroniques mal modelises, perception de
liaisons imparfaite pour les radicaux). Un test annexe avec un filtre
Huckel $\pi$ (gap HOMO--LUMO) a donne 55~\% d'accuracy, soit le taux de
base : aucun pouvoir discriminatif. La conclusion operationnelle est
que \emph{seul xTB peut decider} sur les non-benzenoides ; aucun
pre-filtre rapide n'a ete retenu pour la suite. Le rapport \texttt{xtb}
est donc le critere effectif des resultats qui suivent.

\section{Resultats}

Le Tableau~\ref{tab:pb1-vs-adj57} compare la nouvelle configuration
\texttt{pb1\_adj57} aux deux baselines historiques (\texttt{pb1} de v2,
qui constitue le meilleur reglage anterieur, et la baseline non
contrainte \texttt{baseline\_v2} en regime \texttt{no-freeze, no-table})
sur l'ensemble complet $h_6, h_7, h_8, h_9$, soit
$44 + 126 + 376 + 2\,418 = 2\,964$ squelettes par configuration.

\begin{table}[h]
\centering
\small
\begin{tabular}{lrrrrrr}
\toprule
$h$ & \multicolumn{2}{c}{\texttt{baseline\_v2}} & \multicolumn{2}{c}{\texttt{pb1}} & \multicolumn{2}{c}{\texttt{pb1\_adj57}} \\
\cmidrule(lr){2-3}\cmidrule(lr){4-5}\cmidrule(lr){6-7}
    & \%plan & \#plans & \%plan & \#plans & \%plan & \#plans \\
\midrule
$h_6$ & 70.2 & 896    & 80.5 & 645    & \textbf{88.8} & 293 \\
$h_7$ & 54.8 & 4\,313 & 68.9 & 2\,418 & \textbf{80.8} & 998 \\
$h_8$ & 42.9 & 21\,999 & 61.7 & 9\,391 & \textbf{75.1} & 3\,536 \\
$h_9$ & 21.6 & 121\,692 & 31.3 & 40\,852 & \textbf{39.3} & 13\,580 \\
\bottomrule
\end{tabular}
\caption{Taux de planarite et nombre absolu de solutions planaires par
configuration et par taille de squelette (regime no-freeze, no-table).
\texttt{baseline\_v2} = aucune contrainte additionnelle ; \texttt{pb1} =
contrainte pent-en-bord avec $K_{\text{pb}} = 1$ ; \texttt{pb1\_adj57}
= ajoute la contrainte d'adjacence 5--7 imposant l'appariement local
pent--hept.}
\label{tab:pb1-vs-adj57}
\end{table}

Le Tableau~\ref{tab:gains} resume les gains de \texttt{pb1\_adj57} par
rapport aux deux baselines.

\begin{table}[h]
\centering
\small
\begin{tabular}{lrrrr}
\toprule
$h$ & \%plan baseline\_v2 & \%plan pb1\_adj57 & gain vs pb1 & gain vs baseline \\
\midrule
$h_6$ & 70.2 & 88.8 & $+8.3$  & $+18.6$ \\
$h_7$ & 54.8 & 80.8 & $+11.9$ & $+26.0$ \\
$h_8$ & 42.9 & 75.1 & $+13.4$ & $+32.2$ \\
$h_9$ & 21.6 & 39.3 & $+8.0$  & $+17.7$ \\
\bottomrule
\end{tabular}
\caption{Gains absolus de \texttt{pb1\_adj57} en points de taux de
planarite par rapport a \texttt{pb1} et par rapport a la baseline.}
\label{tab:gains}
\end{table}

\section{Discussion}

\paragraph{Le gain de adj-57 est systematique et substantiel.}
\texttt{pb1\_adj57} ameliore \texttt{pb1} de $+8$ a $+14$ points selon
la taille du squelette, sans aucune exception parmi les quatre valeurs
de $h$ testees. Le gain est maximal sur $h_8$ ($+13.4$ points) et reste
de $+8$ points sur $h_9$, le plus difficile. Compare a la baseline non
contrainte, le saut atteint $+26$ a $+32$ points sur les tailles
intermediaires $h_7$ et $h_8$, soit pratiquement un doublement du
rendement.

\paragraph{Interpretation geometrique : appariement Stone--Wales.}
La contrainte \texttt{adj-57} impose que chaque pentagone admette au
moins un heptagone parmi ses voisins immediats et reciproquement. En
termes de Gauss--Bonnet discret, la courbure $+\pi/3$ d'un pentagone
isole ne peut etre compensee qu'a distance par une autre face de
courbure negative ; sans appariement local, cette courbure doit etre
absorbee par une deformation globale (bombement ou ondulation) qui
disqualifie la molecule en planarite. Avec adj-57, chaque pentagone est
\emph{accole} a un heptagone : le quadripole local $5|6|7$ (ou la
configuration plus large $5|7|7|5$ canoniquement appelee defaut
Stone--Wales) garde une courbure cumulee nulle sur un disque graphique
de rayon 1 ou 2. C'est ce qui explique le gain observe : on ne filtre
pas seulement les configurations a forte courbure, on impose en plus
qu'elles s'organisent localement de maniere a annuler cette courbure.

\paragraph{Trade-off volume / qualite, suite.}
\texttt{pb1\_adj57} produit environ $3$ a $4$ fois moins de solutions
que \texttt{pb1} (par exemple sur $h_7$ : 1\,235 sols vs 3\,510, sur
$h_8$ : 4\,706 sols vs 15\,230). Mais la proportion planaire est telle
que le nombre absolu de molecules planaires par job xTB est environ
multiplie par 1.3 a 1.5. Pour un usage aval avec validation couteuse,
\texttt{pb1\_adj57} domine donc \texttt{pb1} a la fois en rendement
qualitatif et en cout total de validation. Le volume total de plans
produits par \texttt{pb1\_adj57} sur l'union $h_6 \cup h_7 \cup h_8
\cup h_9$ s'eleve a environ 18\,000 molecules, soit un volume
operationnel adequat tout en garantissant une purete superieure a
75~\% jusqu'a $h_8$.

\paragraph{Lecon negative sur les pre-filtres.}
Les deux pre-filtres rapides post-CSP testes (MMFF94 mecanique et
Huckel electronique) ont echoue a discriminer plan / non plan sur les
geometries d'entree brutes. Le mecanisme d'echec est different dans les
deux cas : MMFF94 \emph{aplatit a tort} les configurations
anti-aromatiques (precision reelle 58~\% au lieu des 80~\% attendus),
tandis que Huckel \emph{ne separe pas du tout} les distributions de
gap HOMO--LUMO (55~\% d'accuracy, indistinguable du taux de base). La
planarite d'un non-benzenoide ne se reduit donc ni a un equilibre
purement mecanique, ni a un signal purement electronique : c'est le
couplage geometrie-electronique que xTB capture, et que les modeles
intermediaires bon marche manquent. Cette observation oriente toute
exploration future de pre-filtre vers des methodes hybrides (par
exemple un modele appris sur les paires geometrie / energie produites
par xTB), et non vers les approximations standard mecanique-only ou
$\pi$-only.

\section{Conclusion}

\texttt{experiments\_v3} etablit que la combinaison \texttt{pb1 +
adj-57}, c'est-a-dire le couplage entre la contrainte combinatoire
\emph{au plus un pentagone en bord} et la contrainte structurelle
\emph{appariement pent--hept de premier voisinage}, ameliore de maniere
systematique le taux de planarite des molecules generees par rapport au
meilleur reglage anterieur \texttt{pb1}, avec un gain mesure de $+8$ a
$+14$ points selon la taille du squelette, et un gain cumule de $+18$
a $+32$ points par rapport a la baseline non contrainte. Cette
combinaison est interpretable dans le cadre des defauts topologiques
plans (motif Stone--Wales) comme une condition d'annulation locale de
la courbure de Gauss--Bonnet.

Parallelement, l'echec de deux pre-filtres rapides (MMFF94 et Huckel),
documente quantitativement par une validation par echantillonnage de
5\,996 molecules independantes, identifie une limite operationnelle :
sur les geometries d'entree brutes des non-benzenoides, aucun substitut
bon marche a xTB n'a ete trouve. Toute optimisation future de cout
devra donc soit accepter ce cout xTB, soit construire un modele de
substitution specifiquement entraine sur les donnees du domaine.

\end{document}
